\chapter{Getallen en getallenverzamelingen}\label{ch:g10:numbers}

De wiskunde begint bij de getallen, en niet alle getallen zijn van dezelfde
soort: telgetallen, negatieve getallen, breuken, en getallen als $\sqrt 2$ of
$\pi$ die geen enkele breuk kan uitdrukken. Dit hoofdstuk ordent ze in
geneste families, voert intervallen in om stukken van de getallenlijn te
beschrijven, en gebruikt de \omterm{def:g10:numbers:abs}{absolute waarde} om afstanden tussen getallen te
meten.

\section{De families van getallen}

\begin{definition}[Getallenverzamelingen]\label{def:g10:numbers:sets}
\begin{itemize}
  \item $\N$ is de verzameling van de \emph{natuurlijke getallen}\index{natuurlijk
        getal}: $0, 1, 2, 3, \dots$
  \item $\Z$ is de verzameling van de \emph{gehele getallen}\index{geheel getal}:
        $\dots, -2, -1, 0, 1, 2, \dots$
  \item $\Q$ is de verzameling van de \emph{rationale getallen}\index{rationaal
        getal}: alle quotiënten $\frac{p}{q}$ met $p \in \Z$, $q \in \N$ en
        $q \neq 0$.
  \item $\R$ is de verzameling van de \emph{reële getallen}\index{reëel getal}:
        alle getallen die je op de getallenlijn kunt plaatsen.
\end{itemize}
\end{definition}

\begin{notation}
Het symbool $\in$ lees je als ``behoort tot'': $3 \in \N$, $-5 \in \Z$,
$\frac23 \in \Q$. Het symbool $\subset$ lees je als ``is bevat in'': elk
\omterm{def:g10:numbers:sets}{natuurlijk getal} is een \omterm{def:g10:numbers:sets}{geheel getal}, elk \omterm{def:g10:numbers:sets}{geheel getal} is een \omterm{def:g10:numbers:sets}{rationaal getal}
(bijvoorbeeld $-5 = \frac{-5}{1}$), en elk \omterm{def:g10:numbers:sets}{rationaal getal} is reëel, dus
\[
\N \subset \Z \subset \Q \subset \R .
\]
Een streep door een symbool ontkent het: $\frac12 \notin \Z$.
\end{notation}

\begin{omfigure}
\begin{tikzpicture}[scale=0.9]
  \draw[omExo, thick] (0,0) ellipse (5.4 and 2.9);
  \draw[omProp, thick] (-0.8,0) ellipse (4.0 and 2.2);
  \draw[omThm, thick] (-1.6,0) ellipse (2.6 and 1.5);
  \draw[omDef, thick] (-2.4,0) ellipse (1.3 and 0.8);
  \node[omDef, font=\small\bfseries] at (-2.4,0.45) {$\N$};
  \node[omThm, font=\small\bfseries] at (-1.6,1.1) {$\Z$};
  \node[omProp, font=\small\bfseries] at (-0.8,1.8) {$\Q$};
  \node[omExo, font=\small\bfseries] at (0,2.5) {$\R$};
  \node[font=\small] at (-2.4,-0.15) {$0 \quad 7$};
  \node[font=\small] at (-0.4,-0.9) {$-5$};
  \node[font=\small] at (1.0,-0.4) {$\tfrac23 \quad {-1.7}$};
  \node[font=\small] at (3.9,-0.6) {$\sqrt2 \quad \pi$};
\end{tikzpicture}

{\small De geneste families van getallen: $\N \subset \Z \subset \Q \subset
\R$. Elke ring bevat getallen die niet tot de kleinere behoren.}
\end{omfigure}

\begin{example}\label{ex:g10:numbers:classify}
We plaatsen een paar getallen in de kleinste familie die ze bevat.
\begin{itemize}
  \item $\frac{28}{4} = 7$, dus $\frac{28}{4} \in \N$ hoewel het als breuk
        geschreven staat: vereenvoudig altijd eerst.
  \item $-3.5 = -\frac{35}{10} = -\frac72$ is rationaal maar niet geheel.
  \item $0.333\dots$ (het cijfer $3$ dat eindeloos herhaalt) is gelijk aan
        $\frac13$, een \omterm{def:g10:numbers:sets}{rationaal getal}.
  \item $\sqrt 2$ en $\pi$ zijn reëel maar niet rationaal, zoals we voor
        $\sqrt 2$ meteen zullen zien; zulke getallen heten
        \emph{irrationaal}\index{irrationaal getal}.
\end{itemize}
\end{example}

\begin{theorem}[Irrationaliteit van $\sqrt 2$]\label{thm:g10:numbers:sqrt2}
Het getal $\sqrt 2$ is niet rationaal: geen enkele breuk van \omterm{def:g10:numbers:sets}{gehele getallen}
heeft kwadraat $2$.
\end{theorem}

\begin{proof}
We redeneren uit het ongerijmde, in kleine stappen.
\begin{enumerate}
  \item Stel dat $\sqrt 2 = \frac{p}{q}$ met $p$ en $q$ positieve \omterm{def:g10:numbers:sets}{gehele
        getallen}, en dat de breuk volledig vereenvoudigd is, zodat $p$ en $q$
        niet allebei even zijn.
  \item Kwadrateren van beide leden geeft $2 = \frac{p^2}{q^2}$, dus
        $p^2 = 2q^2$. Bijgevolg is $p^2$ even.
  \item Was $p$ oneven, zeg $p = 2k+1$, dan zou $p^2 = 4k^2 + 4k + 1$ oneven
        zijn. Omdat $p^2$ even is, moet $p$ even zijn: $p = 2k$ voor een zeker
        \omterm{def:g10:numbers:sets}{geheel getal} $k$.
  \item Invullen geeft $(2k)^2 = 2q^2$, dus $4k^2 = 2q^2$, dus $q^2 = 2k^2$.
        Met hetzelfde argument als in stap 3 is $q$ even.
  \item Nu zijn $p$ en $q$ allebei even, in tegenspraak met stap 1. De
        veronderstelling was onmogelijk: $\sqrt 2$ is \omterm{ex:g10:numbers:classify}{irrationaal}.
\end{enumerate}
\end{proof}

\begin{remark}
De \omterm{def:g10:numbers:sets}{rationale getallen} zijn precies de getallen waarvan de decimale
schrijfwijze ofwel afbreekt (zoals $\frac72 = 3.5$) ofwel vanaf een zeker
punt eindeloos hetzelfde blok herhaalt (zoals $\frac13 = 0.333\dots$ of
$\frac{1}{7} = 0.142857\,142857\dots$). Irrationale getallen als
$\sqrt2 = 1.41421356\dots$ hebben decimalen die nooit in een repeterend
patroon vervallen. Die karakterisering nemen we op dit niveau aan.
\end{remark}

\section{De getallenlijn en de intervallen}

De \omterm{def:g10:numbers:sets}{reële getallen} vullen een lijn: kies je een oorsprong $0$ en een
lengte-eenheid, dan hoort bij elk \omterm{def:g10:numbers:sets}{reëel getal} precies één punt.

\begin{definition}[Interval]\label{def:g10:numbers:interval}
Zij $a$ en $b$ \omterm{def:g10:numbers:sets}{reële getallen} met $a < b$. Een
\emph{interval}\index{interval} is de verzameling van alle \omterm{def:g10:numbers:sets}{reële getallen}
tussen twee grenzen. Een vierkante haak betekent dat de grens meetelt, een
ronde haak dat ze wegvalt:
\begin{center}
\begin{tabular}{c|c}
notatie & beschrijving \\
\hline
$\intcc{a}{b}$ & $a \leq x \leq b$ (beide grenzen meegerekend) \\
$\intoo{a}{b}$ & $a < x < b$ (beide grenzen uitgesloten) \\
$\intco{a}{b}$ & $a \leq x < b$ \\
$\intoc{a}{b}$ & $a < x \leq b$ \\
$\intco{a}{+\infty}$ & $x \geq a$ \\
$\intoo{-\infty}{b}$ & $x < b$ \\
\end{tabular}
\end{center}
De symbolen $-\infty$ en $+\infty$ (``oneindig'') zijn geen getallen, alleen
een manier om te zeggen dat het interval eindeloos doorloopt; de haak ernaast
is altijd een ronde haak. De hele lijn $\R$ is het interval
$\intoo{-\infty}{+\infty}$.
\end{definition}

\begin{omfigure}
\begin{tikzpicture}[scale=1.0]
  \draw[-Stealth] (-4.4,0) -- (4.6,0) node[below, font=\small] {$x$};
  \foreach \x in {-4,...,4} \draw (\x,-0.07) -- (\x,0.07);
  \foreach \x in {-4,-2,0,2,4}
    \node[below, font=\small] at (\x,-0.1) {$\x$};
  \draw[omDef, very thick] (-3,0.02) -- (1,0.02);
  \fill[omDef] (-3,0.02) circle (2.2pt);
  \draw[omDef, fill=white, thick] (1,0.02) circle (2.2pt);
  \node[omDef, above, font=\small] at (-1,0.15) {$\intco{-3}{1}$};
  \draw[omThm, very thick] (2,0.02) -- (4.5,0.02);
  \draw[omThm, fill=white, thick] (2,0.02) circle (2.2pt);
  \node[omThm, above, font=\small] at (3.4,0.15) {$\intoo{2}{+\infty}$};
\end{tikzpicture}

{\small Intervallen op de getallenlijn: een volle stip voor een grens die
meetelt, een open stip voor een grens die wegvalt.}
\end{omfigure}

\begin{definition}[Doorsnede en vereniging]\label{def:g10:numbers:interunion}
Zij $I$ en $J$ twee verzamelingen van \omterm{def:g10:numbers:sets}{reële getallen}. De
\emph{doorsnede}\index{doorsnede} $I \cap J$ (``$I$ en $J$'') is de
verzameling van de getallen die tot beide behoren; de
\emph{vereniging}\index{vereniging} $I \cup J$ (``$I$ of $J$'') is de
verzameling van de getallen die tot minstens één van beide behoren.
\end{definition}

\begin{example}\label{ex:g10:numbers:interunion}
Neem $I = \intcc{-1}{3}$ en $J = \intoo{1}{5}$. Teken beide op dezelfde
lijn: ze overlappen tussen $1$ en $3$. Bijgevolg is
\[
I \cap J = \intoc{1}{3},
\qquad
I \cup J = \intco{-1}{5}.
\]
Let op de haken: $1 \notin J$, dus $1 \notin I \cap J$; maar $3 \in I$ en
$3 \in J$, dus $3 \in I \cap J$.
\end{example}

\begin{method}[Rekenen met intervallen]\label{met:g10:numbers:intervals}
Om de \omterm{def:g10:numbers:interunion}{doorsnede} of de \omterm{def:g10:numbers:interunion}{vereniging} van twee intervallen te bepalen:
\begin{enumerate}
  \item teken de getallenlijn en zet de vier grenzen erop;
  \item arceer het eerste \omterm{def:g10:numbers:interval}{interval} boven de lijn en het tweede eronder;
  \item de \omterm{def:g10:numbers:interunion}{doorsnede} is waar de arceringen overlappen, de \omterm{def:g10:numbers:interunion}{vereniging} is waar
        minstens één arcering staat;
  \item beslis elke haak door na te gaan of de grens zelf tot beide
        verzamelingen behoort (\omterm{def:g10:numbers:interunion}{doorsnede}) dan wel tot minstens één
        (\omterm{def:g10:numbers:interunion}{vereniging}).
\end{enumerate}
\end{method}

\section{Absolute waarde en afstand}

\begin{definition}[Absolute waarde]\label{def:g10:numbers:abs}
De \emph{absolute waarde}\index{absolute waarde} van een \omterm{def:g10:numbers:sets}{reëel getal} $x$ is
\[
\abs{x} =
\begin{cases}
x & \text{als } x \geq 0, \\
-x & \text{als } x < 0.
\end{cases}
\]
Zo is $\abs{7} = 7$ en $\abs{-4} = 4$. In alle gevallen is $\abs{x} \geq 0$.
\end{definition}

\begin{proposition}[Afstand op de lijn]\label{prop:g10:numbers:distance}
Voor alle \omterm{def:g10:numbers:sets}{reële getallen} $a$ en $b$ is de afstand tussen de punten $a$ en $b$
op de getallenlijn gelijk aan $\abs{b - a}$. In het bijzonder is $\abs{x}$ de
afstand van $x$ tot $0$.
\end{proposition}

\begin{proof}
Is $b \geq a$, dan is de afstand van $a$ tot $b$ gelijk aan $b - a \geq 0$,
en dat is $\abs{b-a}$. Is $b < a$, dan is de afstand $a - b = -(b - a) > 0$,
en dat is opnieuw $\abs{b - a}$ volgens de definitie van de \omterm{def:g10:numbers:abs}{absolute waarde}.
\end{proof}

\begin{proposition}[Absolute waarde en intervallen]\label{prop:g10:numbers:absinterval}
Zij $a$ een \omterm{def:g10:numbers:sets}{reëel getal} en $r > 0$. Dan geldt
\[
\abs{x - a} \leq r
\quad\text{precies dan als}\quad
x \in \intcc{a - r}{a + r}.
\]
\end{proposition}

\begin{proof}
$\abs{x-a} \leq r$ zegt dat de afstand van $x$ tot $a$ hoogstens $r$ is, dus
dat $x$ aan geen van beide kanten verder dan $r$ van $a$ ligt. De getallen
die daaraan voldoen zijn precies die tussen $a - r$ en $a + r$, grenzen
meegerekend.
\end{proof}

\begin{omfigure}
\begin{tikzpicture}[scale=1.1]
  \draw[-Stealth] (-0.8,0) -- (7,0) node[below, font=\small] {$x$};
  \draw[omDef, very thick] (1.5,0.02) -- (5.5,0.02);
  \fill[omDef] (1.5,0.02) circle (2.2pt);
  \fill[omDef] (5.5,0.02) circle (2.2pt);
  \fill (3.5,0) circle (1.6pt);
  \node[below, font=\small] at (3.5,-0.08) {$a$};
  \node[below, font=\small] at (1.5,-0.08) {$a-r$};
  \node[below, font=\small] at (5.5,-0.08) {$a+r$};
  \draw[Stealth-Stealth, omThm] (3.5,0.35) -- (5.5,0.35)
    node[midway, above, font=\small] {$r$};
  \draw[Stealth-Stealth, omThm] (1.5,0.35) -- (3.5,0.35)
    node[midway, above, font=\small] {$r$};
\end{tikzpicture}

{\small De ongelijkheid $\abs{x - a} \leq r$ beschrijft het \omterm{def:g10:numbers:interval}{interval} van de
getallen op afstand hoogstens $r$ van $a$.}
\end{omfigure}

\begin{example}\label{ex:g10:numbers:abs}
Los $\abs{x - 3} \leq 2$ op. De oplossingen zijn de getallen op afstand
hoogstens $2$ van $3$: het \omterm{def:g10:numbers:interval}{interval} $\intcc{1}{5}$. Omgekeerd heeft het
\omterm{def:g10:numbers:interval}{interval} $\intcc{-1}{7}$ als middelpunt $\frac{-1+7}{2} = 3$ en als straal
$\frac{7-(-1)}{2} = 4$, zodat het beschreven wordt door $\abs{x - 3} \leq 4$.
\end{example}

\section{Benaderingen}

Irrationale getallen, en zelfs de meeste breuken, laten zich niet exact met
eindig veel decimalen schrijven; in de praktijk benaderen we ze dus.

\begin{definition}[Benadering op een gegeven nauwkeurigheid]\label{def:g10:numbers:approx}
Een getal $d$ is een \emph{benadering}\index{benadering} van $x$ op
$10^{-n}$ na wanneer $\abs{x - d} \leq 10^{-n}$. Zowel afkappen als afronden
van de decimale schrijfwijze na het $n$-de cijfer levert zo'n benadering.
\end{definition}

\begin{example}
Uit $\pi = 3.14159\,26\dots$: de afkapping $3.141$ en de afronding $3.142$
zijn allebei \omterm{def:g10:numbers:approx}{benaderingen} van $\pi$ op $10^{-3}$ na. De afronding ligt op
afstand kleiner dan $\frac12 \times 10^{-3}$, de afkapping garandeert enkel
$10^{-3}$. Schrijf je $3.141 \leq \pi \leq 3.142$, dan \emph{klem} je $\pi$
tussen twee decimale grenzen.
\end{example}

\section{Oefeningen}

\begin{exercise}[$\star$]\label{exo:g10:numbers:1}
Geef voor elk getal de kleinste van de verzamelingen $\N$, $\Z$, $\Q$, $\R$
waartoe het behoort:
\[
\frac{15}{3}, \qquad -7, \qquad \frac{22}{7}, \qquad \sqrt{9}, \qquad
\sqrt{10}, \qquad -2.4, \qquad 0 .
\]
\end{exercise}

\begin{exercise}[$\star$]\label{exo:g10:numbers:2}
Schrijf elke uitspraak met intervalnotatie en teken ze daarna op een
getallenlijn: (a) $-2 \leq x < 5$; (b) $x > 3$; (c) $x \leq -1$; (d) de
afstand van $x$ tot $2$ is hoogstens $3$.
\end{exercise}

\begin{exercise}[$\star$]\label{exo:g10:numbers:3}
Bereken $I \cap J$ en $I \cup J$ voor
\[
\text{(a) } I = \intcc{-3}{2},\ J = \intco{0}{4};
\qquad
\text{(b) } I = \intoo{-\infty}{1},\ J = \intco{-2}{+\infty}.
\]
\end{exercise}

\begin{exercise}[$\star$]\label{exo:g10:numbers:4}
Bereken zonder rekenmachine:
\[
\abs{-6}, \qquad \abs{4 - 9}, \qquad \abs{-3 - 5}, \qquad
\abs{\sqrt 2 - 1}, \qquad \abs{1 - \sqrt 2}.
\]
\end{exercise}

\begin{exercise}[$\star$]\label{exo:g10:numbers:5}
Los de vergelijkingen en ongelijkheden op en geef de oplossingsverzamelingen:
\[
\abs{x} = 5, \qquad \abs{x - 1} = 3, \qquad \abs{x - 4} \leq 1, \qquad
\abs{x + 2} < 3 .
\]
(Merk op dat $\abs{x+2} = \abs{x - (-2)}$ een afstand tot $-2$ is.)
\end{exercise}

\begin{exercise}[$\star\star$]\label{exo:g10:numbers:6}
Beschrijf elk \omterm{def:g10:numbers:interval}{interval} door een ongelijkheid van de vorm
$\abs{x - a} \leq r$ of $\abs{x - a} < r$:
\[
\intcc{2}{8}, \qquad \intoo{-5}{1}, \qquad \intcc{-7}{-3}.
\]
\end{exercise}

\begin{exercise}[$\star\star$]\label{exo:g10:numbers:7}
Toon aan dat $0.272727\dots$ (het blok $27$ dat eindeloos herhaalt)
rationaal is. (Tip: noem het $x$ en bereken $100x - x$.)
\end{exercise}

\begin{exercise}[$\star\star$]\label{exo:g10:numbers:8}
Waar of niet waar? Verantwoord elk antwoord met een argument of een
tegenvoorbeeld.
\begin{enumerate}
  \item De som van twee \omterm{def:g10:numbers:sets}{gehele getallen} is geheel.
  \item Het quotiënt van twee \omterm{def:g10:numbers:sets}{gehele getallen} is geheel.
  \item De som van twee \omterm{def:g10:numbers:sets}{rationale getallen} is rationaal.
  \item De som van een rationaal en een \omterm{ex:g10:numbers:classify}{irrationaal getal} is \omterm{ex:g10:numbers:classify}{irrationaal}.
\end{enumerate}
\end{exercise}

\begin{exercise}[$\star\star$]\label{exo:g10:numbers:9}
Klem met behulp van $1.414 \leq \sqrt 2 \leq 1.415$ de getallen $2\sqrt2$,
$\sqrt2 + 3$ en $-\sqrt2$ tussen twee decimale grenzen.
\end{exercise}

\begin{exercise}[$\star\star\star$]\label{exo:g10:numbers:10}
Pas het bewijs van \cref{thm:g10:numbers:sqrt2} aan om aan te tonen dat
$\sqrt 3$ \omterm{ex:g10:numbers:classify}{irrationaal} is. (Vervang ``even'' door ``veelvoud van $3$'': ga
eerst na dat $p$ een veelvoud van $3$ is zodra $p^2$ dat is, door de resten
$0$, $1$, $2$ van $p$ bij deling door $3$ te bekijken.)
\end{exercise}

\section{Opgave: tussen elke twee getallen}

\begin{problem}[{Weekendopgave --- rationale en irrationale getallen
verstrengelen zich: elk interval, hoe klein ook, bevat er oneindig veel van
elke soort, en geen enkele meting kan ze ooit uit elkaar houden}]
\label{pb:g10:numbers:1}
De \omterm{def:g10:numbers:sets}{rationale getallen} lijken een menigte (al die breuken!) en de irrationale
getallen exotische uitzonderingen ($\sqrt2$, $\pi$). Deze opgave onthult het
ware beeld: de twee families \emph{verstrengelen} zich zo fijn dat elk
\omterm{def:g10:numbers:interval}{interval} van de getallenlijn, hoe microscopisch ook, er oneindig veel van
elke soort bevat --- met vreemde gevolgen, zoals dit: geen enkele fysische
meting, hoe nauwkeurig ook, kan ooit beslissen of een lengte rationaal is.

\medskip
\noindent\textbf{Deel I --- De vier koninkrijken.}
\begin{enumerate}
  \item Noem voor elk getal de kleinste van de verzamelingen $\N$, $\Z$,
        $\Q$, $\R$ die het bevat (\cref{def:g10:numbers:sets}): $-7$; \
        $\frac{13}{4}$; \ $\sqrt{16}$; \ $0.121212\ldots$ (denk aan de
        weekendopgave over repeterende decimalen uit het onderbouwvolume); \
        $\sqrt8$; \ $\pi$ (neem zijn irrationaliteit aan --- ze wordt pas in
        de universitaire volumes bewezen).
  \item Bewijs dat $\Q$ stabiel is onder optelling en vermenigvuldiging: zijn
        $x = \frac pq$ en $y = \frac rs$ rationaal, schrijf $x + y$ en $xy$
        dan elk als één breuk.
  \item Leid uit het ongerijmde af: (a) de som van een rationaal en een
        \omterm{ex:g10:numbers:classify}{irrationaal getal} is \omterm{ex:g10:numbers:classify}{irrationaal}; (b) het product van een
        \emph{niet-nul} \omterm{def:g10:numbers:sets}{rationaal getal} en een \omterm{ex:g10:numbers:classify}{irrationaal getal} is
        \omterm{ex:g10:numbers:classify}{irrationaal}.
  \item Toon aan dat de verzameling van de irrationale getallen onder
        \emph{geen} van beide bewerkingen stabiel is: geef twee irrationale
        getallen met een rationale som, en twee met een rationaal product.
  \item Plaats $\sqrt2 + \sqrt8$ en $\sqrt2 \times \sqrt8$ in de vier
        koninkrijken (vereenvoudig eerst $\sqrt8$, met de methode uit het
        onderbouwvolume).
\end{enumerate}

\medskip
\noindent\textbf{Deel II --- De \omterm{def:g10:numbers:sets}{rationale getallen} liggen dicht.}
\begin{enumerate}[resume]
  \item Zoek een \omterm{def:g10:numbers:sets}{rationaal getal} strikt tussen $3.47$ en $3.48$; daarna een
        tweede; beschrijf ten slotte hoe je er zo veel kunt maken als je
        wilt (de inzoombeweging uit de weekendopgave over het ontbreken van
        een volgend getal in het onderbouwvolume, nu met een bewijs in
        zicht).
  \item De algemene stelling. Zij $a < b$ twee willekeurige \omterm{def:g10:numbers:sets}{reële getallen},
        met afstand $g = b - a > 0$. Kies $n$ met $10^{-n} < g$ en bekijk de
        veelvouden van $10^{-n}$ (het decimale rooster met stap $10^{-n}$).
        Leg uit waarom minstens één roosterpunt strikt tussen $a$ en $b$
        valt, en besluit: \emph{elk \omterm{def:g10:numbers:interval}{interval} van positieve lengte bevat een
        \omterm{def:g10:numbers:sets}{rationaal getal}}.
  \item Versterk het besluit: elk zulk \omterm{def:g10:numbers:interval}{interval} bevat \emph{oneindig veel}
        \omterm{def:g10:numbers:sets}{rationale getallen}. (Pas vraag 7 opnieuw toe, binnen een kleiner
        \omterm{def:g10:numbers:interval}{interval}.)
  \item Nu de irrationale getallen: neem bij $a < b$ een \omterm{def:g10:numbers:sets}{rationaal getal} $r$
        strikt ertussen (vraag 7) en bekijk de getallen
        $r + \frac{\sqrt2}{10^k}$. Toon met vraag 3 aan dat ze \omterm{ex:g10:numbers:classify}{irrationaal}
        zijn, en dat ze voor $k$ groot genoeg nog steeds in het \omterm{def:g10:numbers:interval}{interval}
        liggen: \emph{elk \omterm{def:g10:numbers:interval}{interval} bevat ook oneindig veel irrationale
        getallen}.
  \item Twee klassieke raadsels beslecht: bestaat er een kleinste positief
        \omterm{def:g10:numbers:sets}{reëel getal}? En een \omterm{def:g10:numbers:sets}{reëel getal} ``vlak na $3$''? Beantwoord beide met
        de stelling van vraag 7, en groet de kinderversie ervan (de
        weekendopgave over het ontbreken van een volgend getal in het
        onderbouwvolume).
\end{enumerate}

\medskip
\noindent\textbf{Deel III --- De \omterm{def:g10:numbers:abs}{absolute waarde}, de meetkunde van $\R$.}
\begin{enumerate}[resume]
  \item Los op, en druk de oplossingen uit als intervallen of \omterm{def:g10:numbers:interunion}{verenigingen}
        (\cref{prop:g10:numbers:absinterval}): $\abs{x - 5} = 2$; \quad
        $\abs{x - 5} < 2$; \quad $\abs{x + 1} \geq 3$.
  \item Een zuiger moet op $80$ mm gemaakt worden, met een tolerantie van
        $0.05$ mm. Schrijf de eis met een \omterm{def:g10:numbers:abs}{absolute waarde}, en daarna als een
        \omterm{def:g10:numbers:interval}{interval}. Twee zuigers meten $79.97$ en $80.06$ mm: welk oordeel
        krijgen ze?
  \item De driehoeksongelijkheid op de lijn:
        $\abs{a + b} \leq \abs a + \abs b$. Ga ze na op $(a, b) = (3, -5)$ en
        $(-2, -7)$, bewijs ze wanneer $a$ en $b$ hetzelfde teken hebben en
        wanneer ze een tegengesteld teken hebben, en zeg precies wanneer de
        gelijkheid optreedt.
  \item Los $\abs{x - 2} = \abs{x + 4}$ op door het te lezen als een
        gelijkheid van afstanden op de lijn. Welk punt uit de weekendopgave
        over de twee spiegels in het onderbouwvolume heb je zojuist berekend,
        één dimensie lager?
  \item Vereenvoudig $\sqrt{x^2}$ --- voorzichtig. Toets je formule op
        $x = 3$ en $x = -3$, en los er daarna $x^2 < 9$ mee op met een
        \omterm{def:g10:numbers:abs}{absolute waarde}.
\end{enumerate}

\medskip
\noindent\textbf{Deel IV --- Wat geen enkele meting kan beslissen.}
\begin{enumerate}[resume]
  \item Een fysicus meet een staaf: $1.41421 \pm 0.00001$ m. Kan zo'n meting
        --- deze of een toekomstige, nauwkeurigere --- ooit \emph{bewijzen}
        dat de lengte van de staaf \omterm{ex:g10:numbers:classify}{irrationaal} is? Of dat ze rationaal is?
        Gebruik de vragen 7 en 9 op het \omterm{def:g10:numbers:interval}{tolerantie-interval}, en leg uit
        waarom alleen een \emph{bewijs} (zoals voor de diagonaal, in de
        weekendopgave over irrationaliteit in het onderbouwvolume) over
        irrationaliteit kan beslissen.
  \item De afkappingen $1.4$, $1.41$, $1.414$, $1.4142,\dots$ van $\sqrt2$
        zijn allemaal rationaal. Wat tonen ze over hoe dicht $\Q$ tegen elk
        \omterm{ex:g10:numbers:classify}{irrationaal getal} aandringt? Formuleer het algemene feit.
  \item Rekenen met intervallen: gemeten waarden $a = 2.5 \pm 0.1$ en
        $b = 1.2 \pm 0.1$. Klem $a + b$ en $a \times b$ tussen zekere
        grenzen. Welke bewerking tast de nauwkeurigheid het sterkst aan?
  \item Een vuistregel van ingenieurs zegt dat de fouten van onafhankelijke
        metingen bij sommen ``optellen'': druk met de
        driehoeksongelijkheid (vraag 13) uit waarom de fout op $a + b$
        hoogstens de som van de fouten is --- de ongelijkheid \emph{is} de
        vuistregel.
  \item Slotstuk, in een korte alinea: breng het beeld van de reële lijn
        samen dat deze opgave en haar voorgangers hebben opgebouwd --- geen
        volgend getal (de weekendopgave daarover in het onderbouwvolume),
        rationaal $=$ repeterende decimalen (de weekendopgave daarover in het
        onderbouwvolume), beide families dicht (vragen 8 en 9), meting voor
        altijd besluiteloos (vraag 16). Eindig met de vooruitblik die dit
        hoofdstuk nog niet kan bewijzen: in een precieze zin zijn er
        \emph{veel meer} irrationale dan \omterm{def:g10:numbers:sets}{rationale getallen} --- de
        universitaire volumes tellen de oneindigheden.
\end{enumerate}
\end{problem}
